% Lesson 10: shared student/instructor board-supported handout.
% Content: supplied lecture_09(1).tex and lecture_10.tex.
% Structure: converted from the standalone Lesson 10 handout to match the
% wrapper/body layout used by Lessons 01--09.
%
% Source-handling decisions:
% - The narrow-resonance derivation retains the standalone handout's local
%   source approximation and interpretation checkpoints.
% - Instructor reveals are supplied only for algebraic blanks and arithmetic
%   table entries directly implied by the public prompt.
% - Qualitative checkpoints remain student annotation space.
% - Editorial/source limitations are documented in instructor_notes.md.

\providecommand{\dd}{\,\mathrm{d}}
\usetikzlibrary{arrows.meta}
\setlength{\abovedisplayskip}{4pt plus 1pt minus 1pt}
\setlength{\belowdisplayskip}{4pt plus 1pt minus 1pt}
\setlength{\abovedisplayshortskip}{3pt plus 1pt}
\setlength{\belowdisplayshortskip}{3pt plus 1pt}
\setlength{\jot}{2pt}
\tcbset{handout base/.append style={before skip=2.5pt,after skip=2.5pt}}

\HandoutSetup
  {NE 630}
  {10}
  {Resonance Absorption}
  {FNRP \S3.4; review Lesson 09}

\begin{document}
\MakeHandoutTitle
\begin{takeawaybox}[Today's objective]
Approximate the epithermal flux spectrum in an infinite, homogeneous moderator--fuel mixture with resonance absorption. Explain \emph{which assumptions} give the narrow-resonance (NR) approximation and \emph{what its flux depression means}.
\end{takeawaybox}

\begin{handoutcolumns}
\HandoutSection{1}{Start from the slowing-down balance}
Use $E'$ for the energy \emph{before} a collision and $E$ for the energy \emph{after} it. For stationary targets and isotropic elastic scattering in the center-of-mass system,
$$
 \alpha_i=\left(\frac{A_i-1}{A_i+1}\right)^2,
 \qquad
 p_i(E'\!\to E)=\frac{1}{(1-\alpha_i)E'}
$$
on $\alpha_i E'\le E\le E'$; the probability density is zero otherwise. Here $A_i$ is the target-to-neutron mass ratio.

\begin{fixedbox}{Recall: the no-absorption spectrum}
For steady, asymptotic slowing down in one scattering nuclide,
$$ \phi(E)=\frac{K}{\Sigma_s(E)E}. $$
For nearly constant scattering, write $\phi_{\mathrm{bg}}(E)=C/E$, with $C$ setting the normalization.
\end{fixedbox}

\begin{boardbox}{Check the integral limits}
A neutron arriving at $E$ must have started with
$$ E\le E'\le \MathBlank[0.75in]{E/\alpha_i}. $$
For hydrogen, $\alpha=0$ and the upper limit is
$\MathBlank[0.50in]{\infty}$.
\end{boardbox}

\HandoutSection{2}{Add fuel and absorption}
Let $m$ denote the moderator and $f$ the fuel. Define the scattering-in contribution from nuclide $i$ by
$$
 I_i(E)=\int_E^{E/\alpha_i}
 \frac{\Sigma_s^i(E')\phi(E')}{(1-\alpha_i)E'}\dd E'.
$$
The epithermal balance becomes
$$
 \boxed{\Sigma_t(E)\phi(E)=I_m(E)+I_f(E)},
$$
where $\Sigma_t=\Sigma_s^m+\Sigma_s^f+\Sigma_a^m+\Sigma_a^f$.
We neglect leakage and direct neutron births \emph{in this energy interval}; neutrons still enter by slowing down from higher energies.

\begin{checkpoint}[Balance check]
Why does the left side contain $\Sigma_t$, but the right side contains only scattering cross sections?
\RuledLines{2}

\end{checkpoint}

\switchcolumn
\RaggedRight
\HandoutSection{3}{Narrow compared with what?}
A sharp-looking cross-section peak is not enough to justify NR. Compare its energy width $\Delta E_{\mathrm{res}}$ with the energy loss allowed in scattering from \emph{each} nuclide.

For a neutron incident near the resonance energy $E_r$,
$$
 \alpha_i E_r\le E_{\mathrm{out}}\le E_r,
 \qquad
 \Delta E_i=(1-\alpha_i)E_r.
$$
Here $\Delta E_i$ is the \textbf{maximum allowed loss}, used as an energy-transfer scale. A collision can still transfer arbitrarily little energy.

\begin{fixedbox}{Regime assumed for the NR model}
The resonance is narrow on the moderator \emph{and} fuel energy-transfer scales:
$$
 \Delta E_{\mathrm{res}}\ll(1-\alpha_m)E_r,
 \qquad
 \Delta E_{\mathrm{res}}\ll(1-\alpha_f)E_r.
$$
These are approximation conditions, not a claim that every collision jumps across the resonance.
\end{fixedbox}

\begin{boardbox}{A scale comparison (hypothetical resonance)}
Take $E_r=10$ eV and $\Delta E_{\mathrm{res}}=0.50$ eV. Use integer mass ratios for this exercise.
\medskip

\renewcommand{\arraystretch}{1.35}
\par\noindent
\begin{tabularx}{\linewidth}{@{}lcc@{}}
\toprule
Target & $1-\alpha_i$ & $\Delta E_i$ (eV)\\
\midrule
$A=1$ & $1$ & \MathBlank[0.58in]{10}\\
$A=238$ & $0.0167$ & \MathBlank[0.58in]{0.167}\\
\bottomrule
\end{tabularx}\par
\medskip

Which target clearly fails the NR scale comparison?
\RuledLines{2}

\end{boardbox}

\HandoutSection{4}{Approximate the scattering source}
To evaluate $I_m$ and $I_f$ without first knowing the detailed flux, use the assumptions in the slides:
\begin{enumerate}
\item Absorption is negligible \emph{between} resonances, so the surrounding flux has the shape $C/E'$.
\item In the scattering-in integrals, replace scattering by constant, smooth background values $\Sigma_{s,0}^i$.
\item The resonance is narrow enough for both contributions that the detailed resonance region can be neglected in these integrals.
\end{enumerate}
\begin{takeawaybox}[The approximation is made on the right side]
Approximate $\Sigma_s^i(E')\phi(E')$ in $I_i(E)$ by $\Sigma_{s,0}^i C/E'$, but \textbf{retain the energy-dependent total} $\Sigma_t(E)$ on the left.
\end{takeawaybox}
\end{handoutcolumns}

\newpage
\begin{handoutcolumns}
\HandoutSection{5}{Derive the NR spectrum}
For either scattering nuclide, insert the background approximation from Section 4:
\begin{align*}
 I_i(E)
 &\approx \frac{C\Sigma_{s,0}^i}{1-\alpha_i}
    \int_E^{E/\alpha_i}\frac{\dd E'}{E'^2}\\[3pt]
 &=\frac{C\Sigma_{s,0}^i}{1-\alpha_i}
    \left(\frac{1}{E}-\MathBlank[0.50in]{\alpha_i/E}\right)\\[3pt]
 &=\MathBlank[1.00in]{C\Sigma_{s,0}^i/E}.
\end{align*}
For hydrogen, take $\alpha_i=0$ and integrate to $\infty$; the same result follows.

\begin{boardbox}{Put the two contributions back into balance}
Define $\Sigma_{s,0}=\Sigma_{s,0}^m+\Sigma_{s,0}^f$. Then
$$
 \Sigma_t(E)\phi_{\mathrm{NR}}(E)
 \approx \MathBlank[1.15in]{C\Sigma_{s,0}/E}.
$$
Divide by $\Sigma_t(E)$:
$$
 \boxed{\phi_{\mathrm{NR}}(E)\approx
 \frac{C\Sigma_{s,0}}{\Sigma_t(E)E}}
 \quad\Longrightarrow\quad
 \boxed{\phi_{\mathrm{NR}}\propto\frac{1}{\Sigma_t E}}.
$$
This is the NR result in the slides (like FNRP 3.29).
\end{boardbox}

\textbf{Keep the constants straight.} $\Sigma_{s,0}$ is the constant background used in the source approximation, \emph{not} the detailed $\Sigma_s(E)$. The slowing-down supply fixes $C$; this local argument determines the shape.

\begin{checkpoint}[Two limiting checks]
With the same $C$, what happens when:
\begin{enumerate}
\item absorption vanishes and $\Sigma_t=\Sigma_{s,0}$?
\item $\Sigma_t(E_r)$ increases while the approximate scattering-in source stays fixed?
\end{enumerate}
\RuledLines{3}

\end{checkpoint}

\begin{takeawaybox}[Where to stop]
When the resonance is not narrow for fuel scattering, the replacement of $I_f$ is not justified by this argument. Return to the integral balance; do not assume the NR formula remains accurate.
\end{takeawaybox}

\switchcolumn
\RaggedRight
\HandoutSection{6}{Interpret the flux and reaction rate}
Compare the NR spectrum with the same background $\phi_{\mathrm{bg}}=C/E$:
$$
 \boxed{\frac{\phi_{\mathrm{NR}}(E)}{\phi_{\mathrm{bg}}(E)}
 \approx\frac{\Sigma_{s,0}}{\Sigma_t(E)}}.
$$
Thus a resonance peak in $\Sigma_t$ produces a depression in the flux relative to its $1/E$ background.

\begin{boardbox}{Sketch the consequence}
Sketch $E\phi_{\mathrm{NR}}/C$ for an isolated absorption resonance. The dashed line is the undisturbed $1/E$ spectrum in these coordinates.
\par\smallskip
\centering
\begin{tikzpicture}[x=0.49cm,y=1.55cm,>=Latex]
  \draw[->,thin] (0,0)--(10.8,0) node[right] {$E$};
  \draw[->,thin] (0,0)--(0,1.43) node[above right] {$E\phi/C$};
  \draw[dashed,KSUGray] (0.2,1)--(10.2,1);
  \node[left] at (0,1) {$1$};
  \draw[KSUGray!50,densely dotted] (5.2,0)--(5.2,1.1);
  \node[below] at (5.2,0) {$E_r$};

\end{tikzpicture}
\par\raggedright\scriptsize Schematic prompt; no evaluated nuclear data are implied.
\end{boardbox}

\textbf{A derived interpretation.} A lower flux at the absorption peak reduces absorption relative to using an undisturbed $C/E$ flux there. This is \emph{energy self-shielding}; no spatial shielding is needed in this homogeneous model.

\begin{checkpoint}[A flux dip does not mean zero absorption]
For this check only, let $\Sigma_t=\Sigma_{s,0}+\Sigma_a$ and define $x=\Sigma_a/\Sigma_{s,0}$. Hold the approximate scattering-in source $I=C\Sigma_{s,0}/E$ fixed. Then
$$
 \frac{\phi_{\mathrm{NR}}}{\phi_{\mathrm{bg}}}=\frac{1}{1+x},
 \qquad
 \frac{R_a}{I}=\frac{x}{1+x},
 \quad R_a=\Sigma_a\phi_{\mathrm{NR}}.
$$
\renewcommand{\arraystretch}{1.28}
\par\noindent
\begin{tabularx}{\linewidth}{@{}c>{\centering\arraybackslash}X>{\centering\arraybackslash}X@{}}
\toprule
$x$ & $\phi_{\mathrm{NR}}/\phi_{\mathrm{bg}}$ & $R_a/I$\\
\midrule
$0$ & $1$ & $0$\\
$1$ & \MathBlank[0.65in]{1/2} & \MathBlank[0.65in]{1/2}\\
$9$ & \MathBlank[0.65in]{1/10} & \MathBlank[0.65in]{9/10}\\
\bottomrule
\end{tabularx}\par
\end{checkpoint}
\end{handoutcolumns}

\begin{takeawaybox}[Return to the neutron-economy question]
Lesson 08 counted production and absorption; Lesson 09 supplied their flux weighting. Lesson 10 shows that \textbf{resonance cross sections also reshape that weighting}. In rates such as $\int\Sigma_a(E)\phi(E)\dd E$, use the depressed spectrum rather than assuming an undisturbed $1/E$ flux through the resonance.
\end{takeawaybox}
\noindent{\scriptsize\color{KSUGray}Reference: FNRP \S3.4; see the lecture-slide references to Eqs. 3.22, 3.28, and 3.29.}

\end{document}
